\item \textbf{¿Por qué no es posible la siguiente situación?}
La distancia perpendicular de una lámpara a partir de un gran espejo plano es el doble de la distancia perpendicular de una persona al espejo. La luz de la lámpara llega a la persona por dos trayectorias: (1) se desplaza hacia el espejo y se refleja desde el espejo a la persona, y (2) viaja directamente a la persona sin reflejarse. La distancia total recorrida por la luz en el primer caso es $3.10$ veces la distancia recorrida por la luz en el segundo caso.